\documentclass[12pt, a4paper]{report}

%----------------------------------------------------------------------------------------
%	PACKAGES AND THEMES
%----------------------------------------------------------------------------------------
\usepackage[utf8]{inputenc}
\usepackage{fontenc}
\usepackage{mathptmx} % Times New Roman-like font for professional look
\usepackage{geometry}
\usepackage{graphicx}
\usepackage{xcolor}
\usepackage{titlesec}
\usepackage{hyperref}
\usepackage{booktabs} % For professional tables
\usepackage{tikz}     % For the beautiful title page graphics
\usetikzlibrary{calc}
\usepackage{fancyhdr}
\usepackage{lipsum}   % For dummy text if needed
\usepackage{amsmath}  % For mathematical formulas

% Geometry Settings
\geometry{
    top=2.5cm,
    bottom=2.5cm,
    left=3cm,   % Extra margin for binding
    right=2.5cm
}

% Color Definitions (Matching University Themes)
\definecolor{cuiblue}{RGB}{0, 51, 102} % Dark Navy Blue
\definecolor{cuigold}{RGB}{218, 165, 32} % Goldenrod accent

% Hyperlink Setup
\hypersetup{
    colorlinks=true,
    linkcolor=cuiblue,
    filecolor=magenta,
    urlcolor=cyan,
    pdftitle={Automated Documentation Generator},
    pdfpagemode=FullScreen,
}

% Chapter Title Formatting
\titleformat{\chapter}[display]
  {\normalfont\bfseries}{}{0pt}{\Huge}

%----------------------------------------------------------------------------------------
%	DOCUMENT BEGINS
%----------------------------------------------------------------------------------------
\begin{document}

%----------------------------------------------------------------------------------------
%	BEAUTIFUL TITLE PAGE DESIGN
%----------------------------------------------------------------------------------------
\begin{titlepage}
    \begin{tikzpicture}[remember picture, overlay]
        % Sidebar Decoration
        \fill[cuiblue] (current page.north west) rectangle ([xshift=3cm]current page.south west);
        \fill[cuigold] ([xshift=2.8cm]current page.north west) rectangle ([xshift=3cm]current page.south west);

        % Corner Decoration
        \node[opacity=0.1, anchor=north east] at (current page.north east) {\includegraphics[width=10cm]{example-image}}; % Placeholder for Logo watermark
    \end{tikzpicture}

    \centering
    \vspace*{1cm}

    % Header
    {\Large \scshape CECOS University Peshawar}\\[0.5cm]
    {\large \scshape Department of Software Engineering}\\[2cm]

    % Main Title with Rules
    \rule{0.9\textwidth}{1pt} \\[0.4cm]
    { \huge \bfseries \color{cuiblue}   \\[0.3cm] Automated Documentation Generator } \\[0.4cm]
    \rule{0.9\textwidth}{1pt} \\[2cm]

    % Assignment Info
    {\Large \textbf{Subject:} Software Architecture \& Design}\\[0.5cm]
    {\Large \textbf{Project }}\\[2cm]

    % Student and Instructor Details (Side by Side)
    \noindent
    \begin{minipage}{0.45\textwidth}
        \begin{flushleft} \large
            \emph{Submitted By:}\\
            \textbf{Yaseen Ahmad}\\
            \textbf{ID:} CU-4417-C-2023
        \end{flushleft}
    \end{minipage}%
    \begin{minipage}{0.45\textwidth}
        \begin{flushright} \large
            \emph{Submitted To:}\\
            \textbf{Course Instructor}\\
            Dr. Faryal
        \end{flushright}
    \end{minipage}

    \vfill

    % Date
    {\Large \textbf{Dated:} 22-01-2026}\\[1cm]

    % Footer Logo (Optional)
    % \includegraphics[width=3cm]{university_logo.png}

\end{titlepage}

%----------------------------------------------------------------------------------------
%	TABLE OF CONTENTS
%----------------------------------------------------------------------------------------
\tableofcontents
\newpage

%----------------------------------------------------------------------------------------
%	CHAPTER 1: INTRODUCTION
%----------------------------------------------------------------------------------------
\chapter{Introduction}

\section{Overview of the Study}
The domain of software engineering has historically prioritized functional correctness and execution speed over the maintainability of the accompanying documentation...

\section{Research Objectives}
The primary motivation for this project is to reduce the extra load from the software engineers...
\begin{itemize}
    \item \textbf{Economic Motivation:} The time spent by senior engineers writing documentation is expensive.
    \item \textbf{Quality Assurance:} Inconsistent documentation styles across a large organization lead to confusion.
    \item \textbf{Living Documentation:} The ultimate goal is to transform documentation from a static snapshot into a living entity.
\end{itemize}

\section{Scope of the Solution}
The following table outlines the boundaries of the proposed system:

\begin{table}[h]
    \centering
    \caption{Scope of the Automated Documentation Solution}
    \begin{tabular}{@{}lp{8cm}@{}}
        \toprule
        \textbf{Category} & \textbf{Details} \\
        \midrule
        Included & Generating inline comments for functions/methods \\
                 & API documentation from code signatures \\
                 & Function/method descriptions \\
                 & Evaluation against human-written documentation \\
        \midrule
        Excluded & Generating complete user manuals \\
                 & Documentation for binary/compiled code \\
                 & Real-time documentation during coding \\
        \bottomrule
    \end{tabular}
\end{table}

%----------------------------------------------------------------------------------------
%	CHAPTER 2: LITERATURE REVIEW
%----------------------------------------------------------------------------------------
\chapter{Literature Review}

\section{Theoretical Framework}
This project operates in the intersection of software engineering and natural language processing (NLP). The core problem is automating the generation of technical documentation...

%----------------------------------------------------------------------------------------
%	CHAPTER 3: RESEARCH METHODOLOGY
%----------------------------------------------------------------------------------------
\chapter{Research Methodology}

\section{Research Design and Approach}
The system architecture supports the following functional requirements (FRs) and non-functional requirements (NFRs)...

\section{Data Analysis Techniques}
We define the Semantic Similarity between generated ($D_g$) and reference ($D_r$) documentation:
\begin{equation}
    similarity(D_g, D_r) = \frac{v_g \cdot v_r}{\|v_g\| \|v_r\|}
\end{equation}
where $v_g, v_r$ are embeddings from models like BERT.

%----------------------------------------------------------------------------------------
%	CHAPTER 4: RESULTS AND DISCUSSION
%----------------------------------------------------------------------------------------
\chapter{Results and Discussion}

\section{Application Architecture}
The application leverages a microservices approach. Below is the comparison of model selection:

\begin{table}[ht]
    \centering
    \caption{Model Selection and Trade-offs}
    \begin{tabular}{lccc}
        \toprule
        \textbf{Feature} & \textbf{LLaMA 3 (70B)} & \textbf{Mistral Large} & \textbf{GPT-4o} \\
        \midrule
        Reasoning Depth & High & Moderate-High & Very High \\
        Context Window & 8k-128k & 32k & 128k \\
        Privacy & Self-Hosted & Self-Hosted & SaaS \\
        Cost & High CAPEX & Moderate CAPEX & High OPEX \\
        \bottomrule
    \end{tabular}
\end{table}

\section{Discussion of Implications}
\textbf{Economic Impact:} Reduces software maintenance costs (estimated at 60-90\% of total software lifecycle cost).

%----------------------------------------------------------------------------------------
%	CHAPTER 5: CONCLUSION
%----------------------------------------------------------------------------------------
\chapter{Conclusion and Recommendations}

\section{Key Findings}
The architecture provides a robust foundation for building, evaluating, and deploying an automated documentation generation system...

\end{document}